\begin{center}
\begin{circuitikz}[european resistors]
 \draw (0,0) to[sV, l={$V = \VO$}] (0,3)
   to[fuse, l={fuse}, i={$I$}] (3,3) -- (8,3);
 \draw (0,0) -- (8,0);
 \draw (4,3) to[R, l={$R_\mathrm{K}$}, *-*] (4,0);
 \draw (6,3) to[R, l={$R_\mathrm{T}$}, *-*] (6,0);
 \draw (8,3) to[R, l={$R_\mathrm{C}$}] (8,0);
\end{circuitikz}
\end{center}

\begin{abcliste}
\abc Each appliance is at the full mains voltage $V$. The fuse is in series with all appliances together, so the whole current flows through it, the sum of the currents through kettle and toaster:
\begin{align}
 I &= \IF \\
   &= \frac{\V}{\Rk}+\frac{\V}{\Rt} \\
   &= \result{\IP{0}{2}}
\end{align}

\abc The coffee machine may draw at most $I_\mathrm{max} - I$. The smaller its resistance, the larger its current, so its resistance must be at least
\begin{align}
 R_\mathrm{C} &= \RcF \\
   &= \frac{\V}{\Imax - \I} \\
   &= \result{\RcP{0}{2}}
\end{align}
Any machine with a smaller resistance draws more current, and the fuse blows.
\end{abcliste}
